\documentclass[11pt]{beamer}

\usetheme{Madrid}
\usecolortheme{default}
\setbeamertemplate{navigation symbols}{}
\setbeamertemplate{footline}[page number]

\title{Introduction to Traceability in Industry 4.0}
\subtitle{Intelligent Systems for Supply Chain Management}
\author{Bikmetov Daniel \and Krupenskiy Roman \\ 6114-100503D}
\date{Samara, 2025}
\institute{Samara State Aerospace University}

\usepackage[utf8]{inputenc}
\usepackage{amsmath}
\usepackage{amssymb}
\usepackage{booktabs}
\usepackage{xcolor}

\begin{document}

% Title Slide
\frame{\titlepage}

% Outline Slide
\begin{frame}
  \frametitle{Presentation Outline}
  \begin{itemize}
    \item Problem and Goal
    \item Core Idea: Three-Layer Architecture
    \item System Architecture: Ontology-Based Modeling
    \item Data Integration: IoT Data Processing
    \item Implementation: Cloud-Based Platform
    \item Key Results: System Capabilities
    \item Key Challenge \& Solution
    \item Conclusion
  \end{itemize}
\end{frame}

% Slide 3: Problem and Goal
\begin{frame}
  \frametitle{Problem and Goal}
  
  \textbf{Why is intelligent traceability needed?}
  
  \begin{itemize}
    \item Industry 4.0 generates massive heterogeneous data (IoT, sensors, CPS)
    \item Supply chains lack real-time transparency and safety assurance
    \item Current systems cannot integrate data from multiple sources
    \item \textbf{Goal:} Create unified, intelligent traceability for product safety
  \end{itemize}
\end{frame}

% Slide 4: Core Idea
\begin{frame}
  \frametitle{Core Idea: Three-Layer Architecture}
  
  \begin{center}
    \begin{tabular}{|c|}
      \hline
      \textbf{Layer 1 -- Semantic} \\
      Ontology-based knowledge representation \\
      \hline
      \textbf{Layer 2 -- Integration} \\
      Unified data aggregation from heterogeneous sources \\
      \hline
      \textbf{Layer 3 -- Cloud} \\
      Remote access and query capabilities \\
      \hline
    \end{tabular}
  \end{center}
\end{frame}

% Slide 5: Ontology-Based Modeling
\begin{frame}
  \frametitle{System Architecture: Ontology-Based Modeling}
  
  Semantic representation of traceability domain
  
  \textbf{Classes:} Product, Process, Location, Timestamp, Quality
  
  \textbf{Relations:} hasStage, hasQuality, hasLocation
  
  \vspace{0.5cm}
  
  Let $O = \{C, R, I\}$ where:
  \begin{itemize}
    \item $C$ = Concepts (Product, Process, Quality, Location, Agent)
    \item $R$ = Relations (hasStage, hasQuality, hasTrace)
    \item $I$ = Instances (specific products, timestamps)
  \end{itemize}
\end{frame}

% Slide 6: Comparison Table
\begin{frame}[fragile]
  \frametitle{System Architecture: Traditional vs. Proposed}
  
  \tiny
  \begin{center}
    \begin{tabular}{|l|c|c|}
      \hline
      \textbf{Aspect} & \textbf{Traditional} & \textbf{Proposed} \\
      \hline
      Data Integration & Manual, error-prone & Automated, ontology-driven \\
      Interoperability & Limited, siloed systems & Full semantic interoperability \\
      Real-time Access & Delayed reports (batch) & Real-time visibility \\
      Quality Assurance & Batch processing & Continuous monitoring \\
      Compliance & Document-based tracking & Automated tracking \\
      Scalability & Fixed capacity infrastructure & Cloud elastic scaling \\
      Error Detection & Post-incident (reactive) & Real-time (proactive) \\
      \hline
    \end{tabular}
  \end{center}
\end{frame}

% Slide 7: Data Integration
\begin{frame}
  \frametitle{Data Integration: IoT Data Processing}
  
  \textbf{Processing Pipeline}
  
  \begin{itemize}
    \item \textbf{Collection:} Real-time data from sensors, smart devices
    \item \textbf{Normalization:} Map heterogeneous formats to ontology
    \item \textbf{Aggregation:} Merge data into unified knowledge base
  \end{itemize}
  
  \vspace{0.5cm}
  \textbf{Data Transformation Formula}
  
  \[D_{\text{transformed}} = f(D_{\text{raw}}, O) = \text{Map}(D_{\text{raw}}) \rightarrow O_{\text{concepts}}\]
\end{frame}

% Slide 8: Cloud Implementation
\begin{frame}
  \frametitle{Implementation: Cloud-Based Platform}
  
  \begin{itemize}
    \item Common knowledge base accessible to all stakeholders
    \item Query interface for trace retrieval and tracking
    \item Real-time synchronization across supply chain
    \item Access control and data privacy mechanisms
  \end{itemize}
  
  \vspace{0.5cm}
  \textbf{Interoperability Function}
  
  \[I(S_1, S_2) = \text{similarity}(\text{schema}_1, \text{schema}_2, O_{\text{mapping}})\]
\end{frame}

% Slide 9: Key Results
\begin{frame}
  \frametitle{Key Results: System Capabilities}
  
  \textbf{Achieved Outcomes}
  \begin{itemize}
    \item Solved interoperability challenges
    \item Automated data integration
    \item Enhanced transparency
  \end{itemize}
  
  \vspace{0.5cm}
  
  \textbf{Practical Benefits}
  \begin{itemize}
    \item Real-time product tracking
    \item Quality assurance automation
    \item Compliance monitoring
  \end{itemize}
\end{frame}

% Slide 10: Challenge and Solution
\begin{frame}
  \frametitle{Key Challenge \& Solution}
  
  \textbf{Challenge: Data Heterogeneity}
  \begin{itemize}
    \item Multiple source formats (JSON, XML, sensor streams)
  \end{itemize}
  
  \vspace{0.5cm}
  
  \textbf{Solution: Semantic Mapping}
  \begin{itemize}
    \item Ontology defines unified semantic model
    \item Transformation rules map source schemas to ontology
    \item Enables interoperable data exchange
  \end{itemize}
\end{frame}

% Slide 11: Conclusion
\begin{frame}
  \frametitle{Thank You!}
  
  \centering
  \Large
  
  Any questions?
  
  \vspace{1cm}
  
  \normalsize
  \textbf{Contact:} bicmetovdaniel@gmail.com
\end{frame}

\end{document}